\documentclass[../../main]{subfiles}

\begin{document}

\section{Subestruturas Algébricas}
\label{sec:matematica:algebra:estruturas-algebricas:subestruturas}

Nesta seção consideramos estruturas algébricas como modelos de uma assinatura
\[
	\Sigma = \{\sigma_i : n_i\}_{i \in I},
\]
e introduzimos a noção de subestrutura como subconjunto compatível com todas as operações da assinatura.

\subsection{Definição geral de subestrutura}

Seja $(A, \Sigma)$ uma estrutura algébrica e seja $B \subseteq A$ um subconjunto não vazio.

Diz-se que $B$ é uma \textbf{subestrutura} de $A$ se, para toda operação $\sigma_i \in \Sigma$, vale:
\[
	\sigma_i^A(b_1,\dots,b_{n_i}) \in B,
	\quad \forall b_1,\dots,b_{n_i} \in B.
\]

Nesse caso, a restrição das operações de $A$ a $B$ induz uma estrutura algébrica $(B, \Sigma)$.

\subsection{Interpretação estrutural}

Uma subestrutura é um subconjunto fechado sob todas as operações da assinatura.

Em outras palavras, a estrutura algébrica é preservada ao restringir o domínio das operações.

\subsection{Subgrupos}

Se $(G,\ast)$ é um grupo, um subconjunto $H \subseteq G$ é um \textbf{subgrupo} se:

\begin{itemize}
	\item $H$ é não vazio;
	\item para todo $a,b \in H$, tem-se $a \ast b \in H$;
	\item para todo $a \in H$, tem-se $a^{-1} \in H$.
\end{itemize}

Equivalente: $(H,\ast)$ é um grupo com a operação herdada de $G$.

\subsection{Subanéis}

Se $(A,+,\cdot)$ é um anel, um subconjunto $B \subseteq A$ é um \textbf{subanel} se:

\begin{itemize}
	\item $B$ é não vazio;
	\item para todo $a,b \in B$, tem-se $a-b \in B$;
	\item para todo $a,b \in B$, tem-se $a \cdot b \in B$.
\end{itemize}

Assim, $(B,+,\cdot)$ é um anel com operações restritas.

\subsection{Subcorpos}

Um subconjunto $K \subseteq F$ de um corpo $F$ é um \textbf{subcorpo} se:

\begin{itemize}
	\item $K$ contém $0$ e $1$;
	\item é fechado sob adição e multiplicação;
	\item todo elemento não nulo de $K$ possui inverso em $K$.
\end{itemize}

\subsection{Subestruturas como restrição de assinaturas}

Em geral, uma subestrutura preserva a mesma assinatura $\Sigma$, mas restringe o domínio das operações:

\[
	\sigma_i^B = \sigma_i^A|_{B^{n_i}}.
\]

Isso garante que todas as identidades válidas em $A$ continuam válidas em $B$.

\subsection{Propriedade fundamental}

Se $B$ é uma subestrutura de $A$, então toda identidade algébrica satisfeita em $A$ também é satisfeita em $B$.

\subsection{Observação conceitual}

Subestruturas permitem analisar partes internas de uma estrutura algébrica sem perder sua coerência operacional.

Essa noção será essencial para a construção de quocientes, bem como para a análise de invariantes em estruturas algébricas.

\end{document}
